% !TeX root = ../main.tex

\chapter{Agentic AI e Design MCP}
\label{cap:agentic_operations_mcp}

\section{Il Limite dell'Automazione Deterministica}
\label{sec:limite_automazione}

Tradizionalmente l'assegnazione degli ordini sarebbe gestita con un approccio deterministico basato su alberi decisionali (\textit{scripting if-else}). Sebbene questo metodo sia efficiente per carichi di lavoro lineari, si rivela inadeguato nel dominio affrontato: ogni missione è influenzata dalla combinazione di priorità, peso del carico, batteria residua, tasso di usura del mezzo e distanza dall'HUB. Scrivere regole rigide per coprire ogni combinazione produrrebbe codice fragile e difficilmente manutenibile. Gli agenti superano questo limite con un ragionamento probabilistico che valuta dinamicamente i trade-off operativi e si adatta a contesti imprevisti senza pre-programmazione esplicita.

\section{Architettura MCP}
\label{sec:architettura_mcp}

Per garantire confini sicuri, il sistema adotta il \textit{Model Context Protocol} (MCP) come strato di interposizione tra i loop LLM (in \texttt{logistic\_ai\_brain.py}) e l'infrastruttura. L'LLM non possiede credenziali dirette né permessi di root sui database: comunica via HTTP con il server MCP, che espone due categorie di tool:
\begin{itemize}
    \item \textbf{Observer (read-only):} \texttt{get\_drones\_status} (API K8s), \texttt{get\_drones\_telemetry}, \texttt{get\_pending\_orders} (InfluxDB).
    \item \textbf{Remediation (write):} \texttt{send\_mqtt\_command} (assegnazione missioni) e \texttt{scale\_drone\_deployment} (auto-scaling Pod), entrambi soggetti alle policy del Capitolo~\ref{cap:safety_guardrails}.
\end{itemize}

\section{Health Agent e Logistic Agent}
\label{sec:agenti}

\textbf{HealthAgent.} Supervisore della salute infrastrutturale: analizza il rapporto ordini pendenti/capacità di flotta e decide lo scaling dei Pod per preservare gli SLO. Identifica anche guasti fisici: se un drone è in stato \texttt{MAINTENANCE} ma le coordinate divergono dall'HUB, deduce uno schianto e innesca le procedure di recupero.

\textbf{LogisticAgent.} Gestisce la coda ordini e l'assegnazione delle missioni. Incrocia gli ordini con lo stato dei droni \texttt{IDLE}, ponderando tre fattori: (i) peso del pacco (che riduce velocità e aumenta consumo), (ii) distanza euclidea (con calcolo preventivo dell'autonomia), (iii) batteria residua e usura. Ordini \textit{high priority} con peso prossimo a 5.0\,kg non vengono assegnati a droni con usura $>70\%$ o batteria insufficiente. Selezionata la coppia ottimale, l'agente invoca \texttt{send\_mqtt\_command}.

\section{Orchestrazione Dinamica}
\label{sec:logistic_ai_brain}

L'orchestratore (\texttt{logistic\_ai\_brain.py}) coordina gli agenti tramite due meccanismi chiave:
\begin{enumerate}
    \item \textbf{Triage Manager e Context Injection:} l'orchestratore analizza lo snapshot del sistema, decide quale agente attivare e gli inietta nel prompt iniziale solo i dati rilevanti (es. solo i droni \texttt{IDLE} al \texttt{LogisticAgent}), portando l'LLM direttamente alla fase di invocazione dei tool.
    \item \textbf{Esecuzione parallela:} se servono sia manutenzione/scaling sia assegnazione ordini, i due agenti vengono lanciati in parallelo su \texttt{threading.Thread} separati, dimezzando la latenza del ciclo decisionale.
\end{enumerate}
